\section{$x^{n}$-equivalence}

\begin{definition}
\label{def.fps.xneq}
\lean{PowerSeries.XnEquiv}
\leantarget
\leanok
Let $n\in\mathbb{N}$. Let $f,g\in K\left[\left[x\right]\right]$ be two FPSs.
We write $f\overset{x^{n}}{\equiv}g$ if and only if
\[
\text{each }m\in\left\{0,1,\ldots,n\right\}\text{ satisfies }
\left[x^{m}\right]f=\left[x^{m}\right]g.
\]

Thus, we have defined a binary relation $\overset{x^{n}}{\equiv}$ on the set
$K\left[\left[x\right]\right]$. We say that an FPS $f$ is $x^{n}$\emph{-equivalent}
to an FPS $g$ if and only if $f\overset{x^{n}}{\equiv}g$.
\end{definition}

\begin{lemma}
\label{lem.fps.xneq.of-le}
\lean{PowerSeries.XnEquiv.of_le}
\leanhelper
\leanok
Let $m,n\in\mathbb{N}$ with $m\leq n$. If $f\overset{x^{n}}{\equiv}g$,
then $f\overset{x^{m}}{\equiv}g$.
\end{lemma}

\begin{proof}
\leanok
If each $k\in\{0,1,\ldots,n\}$ satisfies $[x^k]f=[x^k]g$, then in particular
each $k\in\{0,1,\ldots,m\}\subseteq\{0,1,\ldots,n\}$ satisfies $[x^k]f=[x^k]g$.
\end{proof}

\begin{lemma}
\label{lem.fps.xneq.zero-iff}
\lean{PowerSeries.XnEquiv_zero_iff}
\leanhelper
\leanok
Two FPSs $f,g\in K[[x]]$ satisfy $f\overset{x^{0}}{\equiv}g$ if and only if
$[x^0]f=[x^0]g$.
\end{lemma}

\begin{proof}
\leanok
This is immediate from the definition, since $\{0,1,\ldots,0\}=\{0\}$.
\end{proof}

\subsection{Helpers for the equivalence relation}

\begin{lemma}
\label{lem.fps.xneq.refl}
\lean{PowerSeries.XnEquiv.refl}
\leanhelper
\leanok
Let $n\in\mathbb{N}$. For any FPS $f\in K[[x]]$, we have
$f\overset{x^{n}}{\equiv}f$.
\end{lemma}

\begin{proof}
\leanok
Immediate from the definition: $[x^m]f=[x^m]f$ for all $m$.
\end{proof}

\begin{lemma}
\label{lem.fps.xneq.symm}
\lean{PowerSeries.XnEquiv.symm}
\leanhelper
\leanok
Let $n\in\mathbb{N}$. If $f,g\in K[[x]]$ satisfy $f\overset{x^{n}}{\equiv}g$,
then $g\overset{x^{n}}{\equiv}f$.
\end{lemma}

\begin{proof}
\leanok
If $[x^m]f=[x^m]g$ for all $m\leq n$, then $[x^m]g=[x^m]f$ by symmetry of equality.
\end{proof}

\begin{lemma}
\label{lem.fps.xneq.trans}
\lean{PowerSeries.XnEquiv.trans}
\leanhelper
\leanok
Let $n\in\mathbb{N}$. If $f,g,h\in K[[x]]$ satisfy $f\overset{x^{n}}{\equiv}g$
and $g\overset{x^{n}}{\equiv}h$, then $f\overset{x^{n}}{\equiv}h$.
\end{lemma}

\begin{proof}
\leanok
If $[x^m]f=[x^m]g$ and $[x^m]g=[x^m]h$ for all $m\leq n$, then
$[x^m]f=[x^m]h$ by transitivity of equality.
\end{proof}

\begin{theorem}
\label{thm.fps.xneq.props}
\lean{PowerSeries.XnEquiv.equivalence, PowerSeries.XnEquiv.add, PowerSeries.XnEquiv.sub, PowerSeries.XnEquiv.mul, PowerSeries.XnEquiv.smul, PowerSeries.XnEquiv.invOfUnit, PowerSeries.XnEquiv.inv, PowerSeries.XnEquiv.div, PowerSeries.XnEquiv.sum, PowerSeries.XnEquiv.prod}
\leantarget
\leanok
Let $n\in\mathbb{N}$. \medskip

\textbf{(a)} The relation $\overset{x^{n}}{\equiv}$ on $K\left[\left[
x\right]\right]$ is an equivalence relation. In other words:

\begin{itemize}
\item This relation is reflexive (i.e., we have $f\overset{x^{n}}{\equiv}f$
for each $f\in K\left[\left[x\right]\right]$).

\item This relation is transitive (i.e., if three FPSs $f,g,h\in K\left[
\left[x\right]\right]$ satisfy $f\overset{x^{n}}{\equiv}g$ and
$g\overset{x^{n}}{\equiv}h$, then $f\overset{x^{n}}{\equiv}h$).

\item This relation is symmetric (i.e., if two FPSs $f,g\in K\left[\left[
x\right]\right]$ satisfy $f\overset{x^{n}}{\equiv}g$, then
$g\overset{x^{n}}{\equiv}f$).
\end{itemize}

\textbf{(b)} If $a,b,c,d\in K\left[\left[x\right]\right]$ are four
FPSs satisfying $a\overset{x^{n}}{\equiv}b$ and $c\overset{x^{n}}{\equiv}d$,
then we also have
\begin{align}
&  a+c\overset{x^{n}}{\equiv}b+d;\label{eq.thm.fps.xneq.props.b.+}\\
&  a-c\overset{x^{n}}{\equiv}b-d;\label{eq.thm.fps.xneq.props.b.-}\\
&  ac\overset{x^{n}}{\equiv}bd. \label{eq.thm.fps.xneq.props.b.*}
\end{align}

\textbf{(c)} If $a,b\in K\left[\left[x\right]\right]$ are two FPSs
satisfying $a\overset{x^{n}}{\equiv}b$, then $\lambda a\overset{x^{n}}{\equiv
}\lambda b$ for each $\lambda\in K$. \medskip

\textbf{(d)} If $a,b\in K\left[\left[x\right]\right]$ are two
invertible FPSs satisfying $a\overset{x^{n}}{\equiv}b$, then $a^{-1}
\overset{x^{n}}{\equiv}b^{-1}$. \medskip

\textbf{(e)} If $a,b,c,d\in K\left[\left[x\right]\right]$ are four
FPSs satisfying $a\overset{x^{n}}{\equiv}b$ and $c\overset{x^{n}}{\equiv}d$,
and if the FPSs $c$ and $d$ are invertible, then we also have
\begin{equation}
\dfrac{a}{c}\overset{x^{n}}{\equiv}\dfrac{b}{d}.
\label{eq.thm.fps.xneq.props.e./}
\end{equation}

\textbf{(f)} Let $S$ be a finite set. Let $\left(a_{s}\right)_{s\in S}\in
K\left[\left[x\right]\right]^{S}$ and $\left(b_{s}\right)_{s\in
S}\in K\left[\left[x\right]\right]^{S}$ be two families of FPSs such
that
\begin{equation}
\text{each }s\in S\text{ satisfies }a_{s}\overset{x^{n}}{\equiv}b_{s}.
\label{eq.thm.fps.xneq.props.e.ass}
\end{equation}
Then, we have
\begin{align}
&  \sum_{s\in S}a_{s}\overset{x^{n}}{\equiv}\sum_{s\in S}b_{s}
;\label{eq.thm.fps.xneq.props.e.+}\\
&  \prod_{s\in S}a_{s}\overset{x^{n}}{\equiv}\prod_{s\in S}b_{s}.
\label{eq.thm.fps.xneq.props.e.*}
\end{align}
\end{theorem}

\begin{proof}
All of these properties are analogous to familiar properties of integer
congruences, except for Theorem \ref{thm.fps.xneq.props} \textbf{(d)},
which is moot for integers (since there are not many integers that are
invertible in $\mathbb{Z}$).

\textbf{(a)} Reflexivity, symmetry and transitivity follow directly from the
corresponding properties of equality of coefficients.

\textbf{(b)} For addition: If $a\overset{x^{n}}{\equiv}b$ and
$c\overset{x^{n}}{\equiv}d$, then for each $m\in\{0,1,\ldots,n\}$ we have
$[x^m](a+c)=[x^m]a+[x^m]c=[x^m]b+[x^m]d=[x^m](b+d)$. Subtraction is
analogous.

For multiplication: We have
$[x^m](ac)=\sum_{i=0}^{m}[x^i]a\cdot[x^{m-i}]c$. Since $m\leq n$, each
$i\leq m$ satisfies $i\leq n$ and $m-i\leq n$, so $[x^i]a=[x^i]b$ and
$[x^{m-i}]c=[x^{m-i}]d$. Thus $[x^m](ac)=[x^m](bd)$.

\textbf{(c)} For each $m\leq n$, we have
$[x^m](\lambda a)=\lambda\cdot[x^m]a=\lambda\cdot[x^m]b=[x^m](\lambda b)$.

\textbf{(d)} Let $a,b\in K[[x]]$ be two invertible FPSs satisfying
$a\overset{x^{n}}{\equiv}b$. We want to show that
$a^{-1}\overset{x^{n}}{\equiv}b^{-1}$.

The FPS $a$ is invertible; thus, its constant term $[x^{0}]a$ is invertible
in $K$ (by Proposition \ref{prop.fps.invertible}).

We prove that each $m\in\{0,1,\ldots,n\}$ satisfies
$[x^{m}](a^{-1})=[x^{m}](b^{-1})$ by strong induction on $m$. Fix some
$k\in\{0,1,\ldots,n\}$, and assume (as induction hypothesis) that
$[x^{m}](a^{-1})=[x^{m}](b^{-1})$ for each $m\in\{0,1,\ldots,k-1\}$.

We know that
\begin{align*}
[x^{k}](aa^{-1})  &  =\sum_{i=0}^{k}[x^{i}]a\cdot[x^{k-i}](a^{-1})\\
&  =[x^{0}]a\cdot[x^{k}](a^{-1})+\sum_{i=1}^{k}[x^{i}]a\cdot[x^{k-i}](a^{-1})
\end{align*}
(here, we have split off the addend for $i=0$ from the sum). Thus,
\[
[x^{0}]a\cdot[x^{k}](a^{-1})+\sum_{i=1}^{k}[x^{i}]a\cdot
[x^{k-i}](a^{-1})=[x^{k}]\underbrace{(aa^{-1})}_{=1}=[x^{k}]1.
\]
We can solve this equation for $[x^{k}](a^{-1})$ (since $[x^{0}]a$ is
invertible), and thus obtain
\[
[x^{k}](a^{-1})=\dfrac{1}{[x^{0}]a}\cdot\left([x^{k}]1-\sum_{i=1}^{k}
[x^{i}]a\cdot[x^{k-i}](a^{-1})\right).
\]
The same argument (applied to $b$ instead of $a$) yields
\[
[x^{k}](b^{-1})=\dfrac{1}{[x^{0}]b}\cdot\left([x^{k}]1-\sum_{i=1}^{k}
[x^{i}]b\cdot[x^{k-i}](b^{-1})\right).
\]
The right hand sides of the latter two equalities are equal (since each
$i\in\{1,2,\ldots,k\}$ satisfies $[x^{i}]a=[x^{i}]b$ as a consequence of
$a\overset{x^{n}}{\equiv}b$, and satisfies $[x^{k-i}](a^{-1})=[x^{k-i}](b^{-1})$
as a consequence of the induction hypothesis, and since we have
$[x^{0}]a=[x^{0}]b$ as a consequence of $a\overset{x^{n}}{\equiv}b$).
Hence, the left hand sides must also be equal. In other words,
$[x^{k}](a^{-1})=[x^{k}](b^{-1})$, which is precisely what we wanted to
prove. Thus, the induction step is complete, so that
$a^{-1}\overset{x^{n}}{\equiv}b^{-1}$ is proved.

\textbf{(e)} Combine parts \textbf{(b)} and \textbf{(d)}: we have
$a\overset{x^{n}}{\equiv}b$ and $c^{-1}\overset{x^{n}}{\equiv}d^{-1}$ (by
\textbf{(d)}), so $ac^{-1}\overset{x^{n}}{\equiv}bd^{-1}$ (by \textbf{(b)}).

\textbf{(f)} The sum case follows by induction on $|S|$ using part
\textbf{(b)} for the induction step (adding one element at a time). The
product case is analogous, using the multiplication part of \textbf{(b)}.
\end{proof}

\begin{lemma}
\label{lem.fps.xneq.neg}
\lean{PowerSeries.XnEquiv.neg}
\leanhelper
\leanok
Let $n\in\mathbb{N}$. If $a,b\in K[[x]]$ satisfy $a\overset{x^{n}}{\equiv}b$,
then $-a\overset{x^{n}}{\equiv}-b$.
\end{lemma}

\begin{proof}
\leanok
For each $m\leq n$, we have
$[x^m](-a)=-[x^m]a=-[x^m]b=[x^m](-b)$.
\end{proof}

\subsection{Helpers for inversion and division}

\begin{lemma}
\label{lem.fps.xneq.invOfUnit-prime}
\lean{PowerSeries.XnEquiv.invOfUnit'}
\leanhelper
\leanok
Let $n\in\mathbb{N}$. Let $f,g\in K[[x]]$ be two FPSs satisfying
$f\overset{x^{n}}{\equiv}g$, and let $u,v$ be units of $K$ satisfying
$u = v$. Then the FPS obtained by inverting $f$ using the unit $u$
is $x^n$-equivalent to the FPS obtained by inverting $g$ using the unit $v$.
\end{lemma}

\begin{proof}
\leanok
The proof proceeds by strong induction on the coefficient index, analogous
to the proof of Theorem~\ref{thm.fps.xneq.props} \textbf{(d)}, using
the equality of unit values $u=v$ and the $x^n$-equivalence of $f$ and $g$.
\end{proof}

\begin{lemma}
\label{lem.fps.xneq.div}
\lean{PowerSeries.XnEquiv.div}
\leanhelper
\leanok
Let $n\in\mathbb{N}$. Let $a,b,c,d\in K[[x]]$ be four FPSs satisfying
$a\overset{x^{n}}{\equiv}b$ and $c\overset{x^{n}}{\equiv}d$,
and let $c$ and $d$ have invertible constant terms. Then
$a \cdot c^{-1}\overset{x^{n}}{\equiv}b \cdot d^{-1}$.
\end{lemma}

\begin{proof}
\leanok
The constant terms of $c$ and $d$ agree (since $c\overset{x^{n}}{\equiv}d$),
so Lemma~\ref{lem.fps.xneq.invOfUnit-prime} gives
$c^{-1}\overset{x^{n}}{\equiv}d^{-1}$. Then
$a \cdot c^{-1}\overset{x^{n}}{\equiv}b \cdot d^{-1}$
by the multiplication rule of Theorem~\ref{thm.fps.xneq.props} \textbf{(b)}.
\end{proof}

\begin{lemma}
\label{lem.fps.xneq.pow}
\lean{PowerSeries.XnEquiv.pow}
\leanhelper
\leanok
Let $n\in\mathbb{N}$. If $a,b\in K[[x]]$ satisfy $a\overset{x^{n}}{\equiv}b$,
then $a^k\overset{x^{n}}{\equiv}b^k$ for every $k\in\mathbb{N}$.
\end{lemma}

\begin{proof}
\leanok
By induction on $k$. The base case $k=0$ gives $a^0=1=b^0$, which is
trivially $x^n$-equivalent to itself. For the induction step, $a^{k+1}=a^k\cdot a
\overset{x^{n}}{\equiv}b^k\cdot b=b^{k+1}$ by the induction hypothesis and
(\ref{eq.thm.fps.xneq.props.b.*}).
\end{proof}

\begin{lemma}
\label{lem.fps.xneq.trunc}
\lean{PowerSeries.xnEquiv_iff_trunc}
\leanhelper
\leanok
Let $n\in\mathbb{N}$. Two FPSs $f,g\in K[[x]]$ satisfy $f\overset{x^{n}}{\equiv}g$
if and only if $\operatorname{trunc}_{n+1}(f)=\operatorname{trunc}_{n+1}(g)$.
\end{lemma}

\begin{proof}
\leanok
The truncation $\operatorname{trunc}_{n+1}(f)$ is the polynomial
$\sum_{k=0}^{n}([x^k]f)\cdot x^k$. Two such truncations are equal if and only
if all their coefficients at degrees $0,1,\ldots,n$ agree, which is exactly the
condition $f\overset{x^{n}}{\equiv}g$.
\end{proof}

\begin{lemma}
\label{lem.fps.xneq.exists-poly}
\lean{PowerSeries.exists_polynomial_xnEquiv}
\leanhelper
\leanok
Let $n\in\mathbb{N}$ and $f\in K[[x]]$. There exists a polynomial
$p\in K[x]$ such that $f\overset{x^{n}}{\equiv}p$.
\end{lemma}

\begin{proof}
\leanok
Take $p=\sum_{k=0}^{n}([x^k]f)\cdot x^k$. Then for each $m\leq n$,
$[x^m]p=[x^m]f$.
\end{proof}

\subsection{Divisibility characterization}

\begin{proposition}
\label{prop.fps.xneq-multiple}
\lean{PowerSeries.xnEquiv_iff_dvd}
\leantarget
\leanok
Let $n\in\mathbb{N}$. Let $f,g\in K\left[\left[x\right]\right]$ be two FPSs.
Then, we have $f\overset{x^{n}}{\equiv}g$ if and only if the FPS $f-g$ is a
multiple of $x^{n+1}$.
\end{proposition}

\begin{proof}
\leanok
A simple consequence of Lemma \ref{lem.fps.muls-of-xn}.

$(\Longrightarrow)$: Assume $f\overset{x^{n}}{\equiv}g$. Then for each
$m<n+1$ (i.e., $m\leq n$), we have $[x^m](f-g)=[x^m]f-[x^m]g=0$.
Thus the first $n+1$ coefficients of $f-g$ are $0$, so $f-g$ is a multiple
of $x^{n+1}$ by Lemma \ref{lem.fps.muls-of-xn}.

$(\Longleftarrow)$: Assume $x^{n+1}\mid(f-g)$. Then for each $m<n+1$, we have
$[x^m](f-g)=0$ (by the power series analogue of Lemma \ref{lem.fps.muls-of-xn}),
so $[x^m]f=[x^m]g$. Since this holds for each $m\leq n$, we obtain
$f\overset{x^{n}}{\equiv}g$.
\end{proof}

\begin{lemma}
\label{lem.fps.xneq-multiple.alt}
\lean{PowerSeries.xnEquiv_iff_sub_eq_mul_X_pow}
\leanhelper
\leanok
Let $n\in\mathbb{N}$. Let $f,g\in K[[x]]$ be two FPSs. Then,
$f\overset{x^{n}}{\equiv}g$ if and only if there exists a $q\in K[[x]]$
such that $f-g=x^{n+1}q$.
\end{lemma}

\begin{proof}
\leanok
This is simply a restatement of Proposition \ref{prop.fps.xneq-multiple},
since ``$f-g$ is a multiple of $x^{n+1}$'' means exactly that there exists
such a $q$.
\end{proof}

\subsection{Composition}

\begin{lemma}
\label{lem.fps.coeff-pow-zero}
\lean{PowerSeries.coeff_pow_eq_zero_of_lt_of_constantCoeff_eq_zero}
\leanhelper
\leanok
Let $f\in K[[x]]$ be an FPS with $[x^0]f=0$. Let $m,k\in\mathbb{N}$ with $m<k$.
Then $[x^m](f^k)=0$.
\end{lemma}

\begin{proof}
\leanok
Since $[x^0]f=0$, the order of $f$ is at least $1$, and thus the order of $f^k$
is at least $k>m$. Hence $[x^m](f^k)=0$.
\end{proof}

\begin{proposition}
\label{prop.fps.xneq.comp}
\lean{PowerSeries.XnEquiv.comp}
\leantarget
\leanok
Let $n\in\mathbb{N}$. Let $a,b,c,d\in K\left[\left[x\right]\right]$ be four
FPSs satisfying $a\overset{x^{n}}{\equiv}b$ and $c\overset{x^{n}}{\equiv}d$
and $[x^{0}]c=0$ and $[x^{0}]d=0$. Then,
\[
a\circ c\overset{x^{n}}{\equiv}b\circ d.
\]
\end{proposition}

\begin{proof}
Write $a$ and $b$ as $a=\sum_{i\in\mathbb{N}}a_{i}x^{i}$ and
$b=\sum_{i\in\mathbb{N}}b_{i}x^{i}$ (with $a_{i},b_{i}\in K$).
Then, $a\overset{x^{n}}{\equiv}b$ means that $a_{i}=b_{i}$ for all $i\leq n$.
Combine this with $c^{i}\overset{x^{n}}{\equiv}d^{i}$ (which holds for all
$i\in\mathbb{N}$ as a consequence of $c\overset{x^{n}}{\equiv}d$ and of
(\ref{eq.thm.fps.xneq.props.e.*})) to obtain the relation
$a_{i}c^{i}\overset{x^{n}}{\equiv}b_{i}d^{i}$ for all $i\leq n$. But this
relation also holds for all $i>n$, since all such $i$ satisfy
$[x^{m}](c^{i})=[x^{m}](d^{i})=0$ for all $m\in\{0,1,\ldots,n\}$ (a
consequence of Lemma \ref{lem.fps.coeff-pow-zero} using $[x^{0}]c=0$ and
$[x^{0}]d=0$). Thus, the relation $a_{i}c^{i}\overset{x^{n}}{\equiv}b_{i}d^{i}$
holds for all $i\in\mathbb{N}$. Summing over all $i$, we find
$a\circ c\overset{x^{n}}{\equiv}b\circ d$.
\end{proof}

\subsection{Equality via $x^n$-equivalence}

\begin{lemma}
\label{lem.fps.eq-iff-forall-xneq}
\lean{PowerSeries.eq_iff_forall_xnEquiv}
\leanhelper
\leanok
Two FPSs $f,g\in K[[x]]$ are equal if and only if $f\overset{x^{n}}{\equiv}g$
for all $n\in\mathbb{N}$.
\end{lemma}

\begin{proof}
\leanok
$(\Longrightarrow)$: If $f=g$, then $f\overset{x^{n}}{\equiv}g$ for every $n$
by reflexivity.

$(\Longleftarrow)$: If $f\overset{x^{n}}{\equiv}g$ for all $n$, then in
particular for each $m\in\mathbb{N}$, taking $n=m$ gives
$[x^m]f=[x^m]g$ (since $m\leq m$). Thus $f=g$ by extensionality
(Lemma~\ref{lem.fps.ext}).
\end{proof}
